%=====================================================================
%  前置部分：扉页、版权与内容简介、译注版说明、前言、目录
%=====================================================================
\begin{titlepage}
  \centering
  \vspace*{2.2cm}
  {\heiti\fontsize{38}{46}\selectfont 量化辛几何导引\par}
  \vspace{0.5cm}
  {\heiti\Large 中文译注版\par}
  \vspace{1.6cm}
  {\large\scshape Introduction to\\[2pt] Quantitative Symplectic Geometry\par}
  \vspace{2.2cm}
  {\large 张俊\quad 祝安铜\quad 编著\par}
  \vfill
  {\large 中国科学技术大学出版社\par}
  \vspace{0.3cm}
  {\small 据中国科学技术大学出版社 2025 年英文版译出\par}
\end{titlepage}

\thispagestyle{empty}
\begin{center}
  {\heiti\large 内容简介}
\end{center}

\noindent
本书系统介绍辛几何与切触几何的核心理论与前沿进展。内容涵盖基础理论、Moser 技巧、哈密顿 Floer 理论和辛同调等；聚焦量子上同调、拉格朗日配边和微局部层论，以代数方法替代传统 Floer 理论；突出定量工具（如持续同调、形不变量等）在刚性研究中的应用；高效证明 Arnold 不动点猜想、Gromov 不可压缩定理等经典结论。

\medskip
\noindent
本书兼顾理论推导与计算实践，每章配备大量例题，既适合入门教学，也可为研究者提供前沿参考。

\vspace{2.2cm}
\noindent
{\heiti 图书在版编目（CIP）数据}

\noindent
量化辛几何导引：英文 / 张俊，祝安铜编著. -- 合肥：中国科学技术大学出版社，2025.6. --（中国科学技术大学一流规划教材）. -- ISBN 978-7-312-06282-7

\noindent
I. O189.3

\noindent
中国国家版本馆 CIP 数据核字第 2025K7N614 号

\vspace{1.6cm}
\noindent
{\heiti 量化辛几何导引（英文版）}\\
LIANGHUA XINJIHE DAOYIN (YINGWEN BAN)

\medskip
\noindent
\begin{tabular}{@{}ll@{}}
出版 & 中国科学技术大学出版社\\
     & 安徽省合肥市金寨路 96 号，230026\\
     & \url{http://press.ustc.edu.cn}\\
     & \url{http://zgkjsdxcbs.tmall.com}\\
印刷 & 安徽省瑞隆印务有限公司\\
发行 & 中国科学技术大学出版社\\
开本 & 787 mm $\times$ 1092 mm\quad 1/16\\
印张 & 23.25\\
字数 & 717 千\\
版次 & 2025 年 6 月第 1 版\\
印次 & 2025 年 6 月第 1 次印刷\\
定价 & 80.00 元
\end{tabular}

\clearpage

%---------------------------------------------------------------------
\chapter*{译注版说明}
\addcontentsline{toc}{chapter}{译注版说明}
\markboth{译注版说明}{译注版说明}

\begin{enumerate}[label=（\arabic*）,leftmargin=3em]
  \item \textbf{底本。}本书据张俊、祝安铜编著《量化辛几何导引》（英文版，中国科学技术大学出版社，2025 年 6 月第 1 版，ISBN 978-7-312-06282-7）的扫描件译出。底本扫描共 363 页，其中正文（第 1--10 章）为印刷页 1--334，对应扫描 PDF 第 9--342 页。

  \item \textbf{缺页说明。}底本扫描件缺失原书附录「Exams」（四次考试，印刷页 335--344，对应 PDF 第 343 页之前约 10 页）。经与目录页码及文献页起始页码核对，可确认该 10 页确系扫描时遗漏，而非原书缺印；本译注版不含该附录，特此说明。

  \item \textbf{编号。}定理、定义、例、注、练习等环境一律沿用原书编号（如「定义 1.1.1」）；公式编号亦照原文手工标注（如式 (1.11)），不依赖自动编号。因此译文与原文的交叉引用（如「由定理 5.4.1」）逐一对得上，方便对照阅读。

  \item \textbf{术语。}采用中文学界通行译名：symplectic geometry 译「辛几何」，contact geometry 译「切触几何」，Lagrangian submanifold 译「拉格朗日子流形」，pseudoholomorphic curve 译「伪全纯曲线」等。重要术语首次出现时以「中文（English）」形式括注；人名一般保留拉丁字母（Gromov、Floer、Arnold 等）。

  \item \textbf{批注。}关键推导处插有 \texttt{\% [Tutor Note: ...]} 形式的 LaTeX 注释，用于提示易错符号约定、作者省略的证明要点、原文疑误及术语依据。注释以 \texttt{\%} 开头，\emph{不进入排版}；如需在 PDF 中显示，可将其改为 \verb|\footnote{}| 或 \verb|\marginpar{}|。全部批注汇总见随附的「审校记录」目录。

  \item \textbf{图表。}原书全部插图均从扫描页按 300 dpi 裁剪为独立 PNG 文件，置于 \texttt{figures/} 目录，并在图注中标明原文图号、印刷页码与 PDF 页码。表格改用 \texttt{booktabs} 重排。

  \item \textbf{编译。}使用 XeLaTeX 编译（文档类为 \texttt{ctexbook}，中文用 Fandol 字体）：
  \begin{verbatim}
  xelatex main.tex
  xelatex main.tex
  \end{verbatim}
  两遍即可生成目录与书签。若教授希望改用 \texttt{article} 体例，只需替换主文件首行文档类，其余环境定义不变。
\end{enumerate}

\clearpage

%---------------------------------------------------------------------
\chapter*{前言}
\addcontentsline{toc}{chapter}{前言}
\markboth{前言}{前言}

本书基于第一作者 2023 年春季与 2024 年春季在中国科学技术大学（USTC）讲授辛几何研究生课程的讲稿。本书是对辛几何与切触几何中标准课题及若干更新课题的入门性介绍。这些内容与许多经典辛几何教材仅部分重叠，我们力求纳入现代机器与数学工具。书中也收录了近期结果，以便迅速把读者带到辛几何与切触几何的研究前沿。

具体而言，本书呈现若干新颖视角。例如，对于 Floer 理论中冗长的介绍——尤其是 Fredholm 型分析——我们将其放入黑箱（并引用若干详尽而漂亮的参考文献），转而采取更定量化的进路，聚焦 Floer 同调的「条形码」结构。由此得到动力学与拓扑中的多种定量应用。特别地，我们给出 Arnold 猜想（哈密顿微分同胚的不动点）与著名的 Gromov 不可压缩定理的快速证明。

再如，在本书后半部分（从第 8 章起），我们几乎不再提及任何 Floer 理论。取而代之，我们详述辛几何中若干其他理论，它们在代数上与 Floer 理论共享某些共同特征：量子上同调（作为哈密顿 Floer（上）同调的代数实现）；Biran--Cornea 的拉格朗日配边理论（作为 Fukaya 范畴的拓扑实现，后者基于拉格朗日 Floer 理论）；微局部层论（作为生成函数理论的量子化理论）。特别地，在第 9 章末，我们给出 Arnold--Givental 猜想的一个证明；该猜想关于哈密顿微分同胚的刚性，体现为零截面与其哈密顿微分同胚像之间的拉格朗日相交性质。该证明完全来自微局部层论，不含任何 $J$-全纯曲线的输入。

本书包含大量习题，几乎全部来自 USTC 课程中的常规作业。此外，附录中还收录了四次考试（两次期中、两次期末）。我们强烈建议读者自行尝试这些习题与考题，因为它们或是对正文某些细节的必要验证，或是由相关语境中的重要结果快速导出的有用结论。

我们感谢许多人。没有学生们在课堂上规律的出席与积极互动，本书无从谈起。我们感谢 USTC 优良的科研与教学环境，它提供了检验所选材料是否适合教学与理解的宝贵机会。课堂内外的众多反馈纠正了初稿中的不准确与错误。特别地，Yichen Yao 作为第一作者两个学期的助教做了出色的工作，我们感谢他的努力与帮助。我们也感谢 Shengxi Jin、Hexiang Wang 和 Matthias Meiwes 阅读本书草稿并提供宝贵反馈。最后，我们感谢 Sungho Kim 博士将本书作为其 USTC 研究生辛几何课程的主要参考书。

我们阐述材料的方式反映了我们在辛几何中的个人思考与工作风格，这在很大程度上受到第一作者博士导师 Michael Usher 教授与博士后导师 Leonid Polterovich 教授的影响与塑造。书中的许多概念与证明借自他们的论文。因此，我们感谢他们对辛几何与切触几何的贡献，使这一学科得以健康地发展繁荣。

本书得到以下资助：2023 年度中国科学技术大学研究生教育创新计划——优秀教材出版项目。第一作者部分受国家重点研发计划 No.~2023YFA1010500、国家自然科学基金 No.~12301081、NSFC No.~12361141812 以及 USTC 双一流建设研究基金资助。

\vspace{1.5em}
\begin{flushright}
张俊\\
2024 年 7 月 16 日
\end{flushright}

\clearpage
\tableofcontents
